\section{Related work}\label{sec:related_work}

RDD integrates several established fields: input minimisation, over-the-air
fuzzing, oracle stabilisation, language models for testing, crash deduplication,
and commercial tools. While each area is well developed individually, we know of no
prior work combining them through adaptive minimisation guided by a statistical
reproduction oracle on a noisy wireless stack.

\subsection{Input minimisation under noisy verdicts}

The primary algorithm, \texttt{ddmin}~\cite{ddmin2002}, is a general-purpose
reducer that identifies a locally minimal failing input. It has been the standard
minimiser in compiler testing and bug reproduction for approximately twenty years;
\textsc{C-Reduce}~\cite{creduce2012} and \textsc{Perses}~\cite{perses2018} are
prominent syntax-guided refinements for structured source programs.
Hierarchical Delta Debugging~\cite{hdd2006} adapts \texttt{ddmin} for
tree-structured inputs, which differs from the noisy verdict problem addressed
here. Probabilistic Delta Debugging (\textsc{ProbDD})~\cite{probdd2021} employs a
Bayesian estimate to determine which elements are significant and to reduce oracle
calls. However, these methods treat the oracle as \emph{deterministic}, making
assumptions about the input rather than the verdict. Probabilistic Monotonicity
Assessment (\textsc{PMA})~\cite{pma2025} departs from this,
modelling the search space's \emph{monotonicity} probabilistically to avoid
redundant tests. \textsc{PMA} is the closest related work; it cites noise and
intermittent failures as motivation, but it accepts each verdict as given and does
not \emph{stabilise} it with repeated trials, sequential tests or an error bound. All these methods assume a
reliable oracle, which wireless environments do not provide, so \texttt{ddmin}'s
1-minimality guarantee does not hold. RDD adopts a different strategy: it maintains
a classical search but prioritises producing a reliable verdict, complementing
rather than competing with these minimisers. Delta debugging has previously been
applied to non-deterministic targets. For example, Pontolillo and
Mousavi~\cite{quantumdd2024} reduce quantum-program regressions using a
fixed-sample property check on a simulator, not a sequential, bounded-error
decision on a real channel.

\subsection{Wireless fuzzing: discovery versus reproduction}

Over-the-air protocol fuzzing is a well-established field.
\textsc{SweynTooth}~\cite{sweyntooth2020} identified eleven BLE vulnerabilities
in chips from seven vendors, and \textsc{BrakTooth}~\cite{braktooth2022} eighteen in
Bluetooth Classic devices from eleven vendors. Both projects conclude at discovery, requiring analysts to
manually construct a clean PoC. \textsc{U-Fuzz}~\cite{ufuzz2024} advances further:
after stateful fuzzing of IoT and 5G protocols, it performs a post-fuzzing step to
identify the minimal set of mutated packets that cause a crash. This makes U-Fuzz
RDD's closest counterpart, but the distinction lies in the oracle. U-Fuzz walks
backwards from the crash replaying one mutated packet at a time in a fresh session
until the crash recurs, states no retry policy, and labels crashes whose
reproduction is unstable as anomalies, while RDD employs a sequential,
bounded-error verdict.
\textsc{AirBugCatcher}~\cite{airbugcatcher2024} is the most direct comparator, as it
also addresses noisy-verdict problems. However, it examines suspect combinations up
to a small fixed size and runs each generated test case once within a trial budget,
rather than stabilising the verdict. In its study, a plain replay baseline,
repeated five times to average out the medium, reproduced only 13 to 23 of 190
crashes on its Bluetooth target and none on the other four~\cite{airbugcatcher2024}. % AirBugCatcher RQ5, Table 7
RDD builds upon these concepts rather than competing with them. Three improvements
have been noted: bounded enumeration is replaced by \texttt{ddmin}, and one attempt
per candidate is replaced by a statistical oracle. The third improvement is in the
identity layer: \textsc{AirBugCatcher} matches crashes by exact source-line trace,
by an address-offset pattern on memory traces, or by packet state when no log
exists, and counts a crash with a non-matching identifier as ``unexpected'';
RDD uses a two-tier cascade tolerant of truncated and garbled reports.

\subsection{Oracle stabilisation and flaky tests}

The oracle problem --- determining whether an outcome is truly a failure --- has
been a focus of testing research for over a decade. Barr et
al.~\cite{oracleproblem2015} review scenarios where no perfect oracle exists, and
Luo et al.~\cite{flakytests2014} discuss flaky tests where verdicts vary across
identical runs. Both works identify the issue but do not provide a decision rule
for noisy \emph{physical} targets. The Sequential Probability Ratio Test
(SPRT)~\cite{sprt1945} addresses this by offering a rule that makes decisions within
a set budget while controlling both error rates. Applying SPRT to the reproduction
oracle introduces Wald's guarantees, reducing false credit from the baseline's
0.632 to 0.000 in all stress tests. Using a statistical test as an oracle is not new: Guderlei and
Mayer~\cite{statoracle2007} test random-output programs by comparing independent
output samples under a metamorphic relation with a fixed-sample hypothesis test, an
offline batch decision, whereas RDD makes sequential, bounded-error verdicts during
the reduction process.

\subsection{Language models (LLMs) as testing oracles}

Language models are now applied in testing in various ways.
\textsc{LPR}~\cite{lpr2024} combines a syntax-guided reducer with LLM-proposed
rewrites, but the model only suggests changes and the oracle remains a
deterministic recompile. \textsc{ChatAFL}~\cite{chatafl2024} allows an LLM to guide
greybox fuzzing of RFC protocols on wired targets, and
\textsc{Candor}~\cite{candor2025} uses multiple models to reach consensus and
generate oracles offline. The most structurally similar approach is
\textsc{DDOR}~\cite{ddor2026}, a recent preprint, which incorporates LLM judging
within a delta-debugging loop for chatbot over-refusal. While similar to RDD, DDOR
targets a replayable software endpoint, without wireless noise or a statistical
reproduction oracle. Another recent effort applies delta debugging to
LLM-integrated systems~\cite{ddllm2026} with semantic markers; it treats the model
as a probabilistic oracle and classifies each subset by a fixed number of repeated
trials with a failure threshold, a majority vote rather than a sequential,
error-bounded test. Ensemble frameworks differ:
self-consistency~\cite{selfconsistency2023} uses a single model's majority vote,
cross-model voting yields only minor gains once voters are
correlated~\cite{condorcetllm2024}, LLM-as-a-judge approaches approximate human
grading~\cite{llmjudge2023}, and unanimous open-weight juries review LLM-written
code offline~\cite{llmjuries2026}. RDD is fundamentally different: its L3 is a
\emph{single} crash-identity judge, used only when the deterministic L2 is uncertain,
similar to FrugalGPT-style escalation~\cite{frugalgpt2023}. The stabiliser is the
SPRT oracle, not an ensemble, and the model is used to infer identity rather than
reproduction.

\subsection{Crash identity and deduplication}

An identity layer is required to match crashes across noisy trials. Beyond exact
stack-hash matching, \emph{Drain}~\cite{drain2017} groups similar logs by mining log
templates online, and Semantic Crash Bucketing~\cite{scb2018} groups crashes based
on whether a simple program change resolves them. Igor~\cite{igor2021} clusters by
root cause using minimised execution traces, offering greater precision, but like
the others, this is an offline triage step. RDD integrates these concepts into a
two-tier identity system: a deterministic L2 that considers stack, site, fault
class, and log templates, and refers only the remaining uncertain cases to an
open-weight model~\cite{llama3_2024}. This cost-gated cascade is incorporated into
the reproduction oracle, unlike previous approaches that treat identity as a
separate triage step. The closest comparison is GPTrace~\cite{gptrace2026}, which
deduplicates crashes using language-model embeddings offline, whereas RDD's judge
determines identity during the process as part of the reproduction verdict.

\subsection{The marketplace}

Commercial wireless fuzzing is well established. Defensics~\cite{defensics} provides
Bluetooth and Wi-Fi suites that transmit malformed traffic over the air, and
Keysight's IoT Security Assessment offers similar capabilities for radio stacks
and incorporated the \textsc{BrakTooth} findings, whose research Keysight
funded~\cite{keysightiot,keysightbraktooth2021,braktooth2022}. Both tools
identify failures but do not perform minimisation. Crash triage is also commercial:
Google's ClusterFuzz~\cite{clusterfuzz} deduplicates and minimises test cases, but
only on locally executed software where results are repeatable. A market review
found no tool that combines adaptive minimisation with a non-deterministic
over-the-air target; these features are offered separately.

\subsection{Where RDD sits}

Input minimisation, wireless discovery, and crash identity are all well-developed
fields, but have not previously been combined. To our knowledge, no prior tool uses
\texttt{ddmin} with a \emph{statistical} reproduction oracle that is SPRT-stabilised
and false-credit-bounded on a non-deterministic wireless target, together with a
tiered identity cascade for the varied dumps caused by noise. Other tools are
similar but not equivalent. \textsc{U-Fuzz} minimises wireless crashes, but only
with a single replay rather than a stabilised verdict. \textsc{PMA} measures
monotonicity but does not provide a bounded-error verdict. \textsc{DDOR} reduces against a replayable
endpoint, and the LLM-system reducer votes over a fixed number of trials. The quantum reducer uses a
fixed-sample check. That is the position RDD takes.
